\begin{statement}[C.4.1; Known, with derivation: Siegel functions and local heights]
For $a=(a_1,a_2)\in\mathbb R^2\setminus\mathbb Z^2$ and
$v=a_1\tau+a_2$, define
\[
 g_a(\tau)=-e^{\pi ia_2(a_1-1)}q^{B_2(a_1)/2}
 (1-e^{2\pi iv})\prod_{r\ge1}
 (1-q^re^{2\pi iv})(1-q^re^{-2\pi iv}),
 \quad B_2(x)=x^2-x+\tfrac16.
\]
Then
\[
 g_a(\tau)=i e^{\pi ia_1v}
       \frac{\vartheta_1(v,\tau)}{\eta(\tau)},\qquad
 N(v,\tau)=|g_a(\tau)|.
\]
When $a\in n^{-1}\mathbb Z^2\setminus\mathbb Z^2$, $n\ge2$,
$g_a^{12n}$ is a modular unit for $\Gamma(n)$ with Fourier coefficients
in $\mathbb Q(\zeta_n)$ in the variable $q^{1/n}$. On every complex
elliptic curve the model-independent Néron local function for $(O)$ is
\[
 \lambda_\infty(Q)=-\log N(Q).
\]
Thus $D_{\rm part}=-(6/\pi)(\lambda_\infty(P)+\log(2\sqrt3))$.
\end{statement}
\begin{proof}
Dividing the two product expressions gives the first identity exactly,
including the factor $i$. Its absolute value is the asserted norm.
This agrees with the convention in \cite{Shin2015}, equation (1).
For completeness the transformation argument needed here can be
checked directly. For $T:\tau\mapsto\tau+1$ and
$S:\tau\mapsto-1/\tau$,
\[
 g_a(T\tau)=e^{\pi i/6}g_{aT}(\tau),\qquad
 g_a(S\tau)=-i g_{aS}(\tau);
\]
the usual row action is used on $a$. For $b=(b_1,b_2)\in\mathbb Z^2$,
\[
 \frac{g_{a+b}}{g_a}
 =\exp\bigl(\pi i(a_1b_2-a_2b_1+b_1b_2+b_1+b_2)\bigr).
\]
These follow from the two Jacobi transformations and the lattice
translations. They are the cases of Proposition 2.3(ii)--(iii) of
\cite{Shin2015} used here. The multipliers on $\Gamma(n)$ are killed
by $12n$. The product gives a Laurent series in $q^{1/n}$ with
coefficients in $\mathbb Q(\zeta_n)$, because
$6nB_2(r/n)\in n^{-1}\mathbb Z$ and its leading unit is an $n$th
root of unity. The transformed products prove meromorphy at every
cusp. No factor has a zero or pole in the upper half-plane. This
proves precisely the modular-unit assertion, also recorded in
Remark 2.2 of that source; its hypotheses $n\ge2$ and a nonintegral
rational pair hold here.

For the local-height normalization, $-U_\Lambda$ has a logarithmic
pole at $O$, mean zero, and Laplacian
$-2\pi\delta_O+2\pi\,d\mu_\Lambda$. This is the unique normalized
Green function for the degree-one divisor $(O)$: two such functions
differ by a harmonic function on the compact curve and their means
remove the constant. Equivalently its Weierstrass expansion is
\[
 -U_\Lambda(z)-\tfrac12\log|x(z)|
             \longrightarrow-\tfrac1{12}\log|\Delta_\Lambda|.
\]
The transformation of $x$ and $\Delta$ under a Weierstrass homothety
cancels in this expression. This is the model-independent local
height normalization, rather than a local function with an arbitrary
additive constant.
\end{proof}
